\documentclass[11pt]{amsart}
\usepackage[margin=1in]{geometry}
\usepackage{amsmath,amssymb,amsthm}
\usepackage{array}
\usepackage{graphicx}
\usepackage[hidelinks]{hyperref}

\graphicspath{{./}{../zenodo_package/}}

\newtheorem{theorem}{Theorem}[section]
\newtheorem{proposition}[theorem]{Proposition}
\newtheorem{corollary}[theorem]{Corollary}
\newtheorem{lemma}[theorem]{Lemma}

\theoremstyle{definition}
\newtheorem{definition}[theorem]{Definition}

\theoremstyle{remark}
\newtheorem*{remark}{Remark}

\newcommand{\Gq}{G_q}
\newcommand{\Gf}{G_5}
\newcommand{\phif}{\varphi_5}
\newcommand{\phiq}{\varphi_q}
\newcommand{\Zf}{Z_{G_5}}
\newcommand{\Zq}{Z_{G_q}}
\newcommand{\Lp}{L_{s,+}}
\newcommand{\Lm}{L_{s,-}}
\newcommand{\detp}{\det(1-L_{s,+})}
\newcommand{\ReC}{\operatorname{Re}}
\newcommand{\ImC}{\operatorname{Im}}

\title[Certified $G_5$ transfer-operator zeros]{Certified zeros of the transfer-operator determinant for the Hecke triangle group $G_5$}

\author{[authors]}

\date{DRAFT --- not for distribution}

\subjclass[2020]{Primary 11F72, 37C30; Secondary 11Y35, 58J50, 30C15}

\keywords{Hecke triangle group, transfer operator, Selberg zeta function, scattering determinant, interval arithmetic, certified computation}

\emergencystretch=2.5em
\begin{document}

\begin{abstract}
The plus-sector transfer-operator determinant $\det(1-L_{s,+})$ of the
Hecke triangle group $G_5$ has two zeros $s_1$ and $s_2$, each of
multiplicity one, confined by interval arithmetic to explicit boxes of
coordinate half-width $10^{-6}$ whose closed real-part intervals are
disjoint. The computation upgrades the numerical zeros recorded by
Str\"omberg for $q=3,4,5$, which that work explicitly described as
non-rigorous, to a certified argument-principle standing: each box is
covered by a closed contour on which a finite-section Fredholm
determinant is enclosed in Arb ball arithmetic with directed rounding,
the enclosure excludes zero, and the certified winding number is $1$.

A published factorization of the Selberg zeta function of $G_5$,
together with a published divisor theorem for cofinite orbifolds,
would promote each certified zero of $\det(1-L_{s,+})$ to a zero of
$Z_{G_5}$ of multiplicity at least one and, after the scattering
functional equation, to a zero of the scalar scattering determinant
$\varphi_5$ in the strip $1/2<\Re s<1$. That promotion is stated as a
conditional proposition, with every joint labelled certified, cited,
or open. No infinitude law for scattering zeros is claimed.
\end{abstract}

\maketitle

\section{Introduction}

For the modular group $\mathrm{PSL}(2,\mathbb Z)$ the scattering
determinant of the cusp is an explicit ratio of completed Riemann zeta
functions, so its nonreal zeros and poles are governed by the zeros of
$\zeta$ and, on the Riemann hypothesis, lie on the critical line
$\Re s=1/2$. The same mechanism operates for every arithmetic Fuchsian
group: number-theoretic $L$-functions supply the scattering data and
pin its zeros to the line. For a non-arithmetic group no such
$L$-function mechanism exists, and the Phillips--Sarnak theory of
cusp-form dissolution predicts that spectral data should instead
migrate into genuinely off-line resonances. Deciding whether the
resonances of a \emph{specific} non-arithmetic group actually leave the
line is then a quantitative question about one meromorphic function,
and the smallest natural test case is the Hecke triangle group $G_5$,
the first non-arithmetic member of the simplest one-parameter family of
Fuchsian groups. Str\"omberg computed numerical candidates for its
lowest resonances and stated explicitly that the computation was not
rigorous. This paper closes that gap for two of them: each is confined
to an explicit box by interval arithmetic with a certified winding
number, and the two boxes have disjoint real parts. The certification
is unconditional; what the two certified zeros do and do not imply
about scattering resonances is stated as a separate, fully labelled
conditional proposition.




Let $q\ge 3$ be an integer. The Hecke triangle group $\Gq$ is the
subgroup of $\mathrm{PSL}(2,\mathbb R)$ generated by the involution
$z\mapsto -1/z$ and the translation $z\mapsto z+\lambda_q$, where
$\lambda_q=2\cos(\pi/q)$. The quotient $X_q=\Gq\backslash\mathbb H$ is
a cofinite hyperbolic orbifold with one cusp and two elliptic points,
of orders $2$ and $q$. For $q=5$ one has $\lambda_5=2\cos(\pi/5)$, and
$G_5$ is non-arithmetic \cite{Takeuchi1977}. Write $\Zq$ for the
Selberg zeta function of $\Gq$, and write $L_{s,\pm}$ for the
plus/minus sectors of the Mayer--M\"uhlenbruch--Str\"omberg transfer
operator of $\Gq$ \cite{MMS2012}. The object certified in this paper
is the plus-sector Fredholm determinant $\detp$.

\begin{definition}
\label{def:resonance}
A \emph{resonance} of $X_q$ is a pole of the meromorphically continued
scalar scattering determinant $\phiq$ of the trivial character. Under
the divisor theorem recalled in Theorem~\ref{thm:fjs} below, every
nonreal zero of $\Zq$ in the strip $0<\Re s<1/2$ is such a pole, with
matching multiplicity, and conversely.
\end{definition}

Str\"omberg \cite{Stromberg2009} computed numerical zeros of
$\Zq$ for $q=3,4,5$ from truncated transfer-operator determinants and
stated that the computation is not rigorous. The present paper
supplies interval certificates for two zeros of $\detp$ at $q=5$.
Winding number $1$ on a contour that excludes zero certifies a simple
zero of $\detp$ in each box. It does not, by itself, certify
simplicity of a zero of $\Zf$: in the factorization of
\cite[Theorem 6.4]{MMS2012},
\begin{equation}
\label{eq:mms-factor}
\Zq(s)=\frac{\det(1-L_{s,+})\,\det(1-L_{s,-})}{\det(1-K_s)},
\end{equation}
the minus-sector factor may vanish at the same point. Accordingly
$\Zf$ is asserted only to vanish with multiplicity at least that of
the plus factor, and only after the joints of
Proposition~\ref{prop:joints} that lie outside the certificates are
granted.

The paper does not prove that every finite $q$ has infinitely many
off-critical-line scattering zeros. That fixed-$q$ growth statement
is discussed only in \S\ref{sec:related}, as an open comparison with
Selberg, Hejhal, and Kelmer.

The contents are as follows. Section~\ref{sec:background} records the
cited analytic apparatus. Section~\ref{sec:method} describes the
interval certificates. Section~\ref{sec:theorem} states the
unconditional two-pin theorem. Section~\ref{sec:joints} enumerates
the thirteen joints from the plus-factor zeros to a scattering
consequence, labelled certified, cited, or open.
Section~\ref{sec:related} compares prior theorems with the present
certificates. Sections~\ref{sec:lean} and~\ref{sec:data} record the
Lean artifacts at their actual strength and the data archive.

\subsection*{Notation}
\begin{center}
\small
\begin{tabular}{@{}>{\raggedright\arraybackslash}p{0.22\textwidth}
>{\raggedright\arraybackslash}p{0.68\textwidth}@{}}
\hline
Symbol & Meaning \\
\hline
$q\ge 3$ & integer elliptic order \\
$\Gq$ & Hecke triangle group $\langle z\mapsto-1/z,\, z\mapsto z+\lambda_q\rangle$ \\
$\lambda_q$ & $2\cos(\pi/q)$ \\
$X_q$ & cofinite orbifold $\Gq\backslash\mathbb H$ \\
$\Zq$ & Selberg zeta function of $\Gq$ \\
$L_{s,\pm}$ & plus/minus Mayer--M\"uhlenbruch--Str\"omberg transfer operators \\
$\detp$ & plus-sector Fredholm determinant \\
$K_s$ & MMS correction operator in the factorization \eqref{eq:mms-factor} \\
$\phiq$ & scalar trivial-character scattering determinant of $X_q$ \\
$\Omega^*$ & $\{\Re s>1/2\}\cup\{\Re s>0,\,\Im s>1\}$ \\
$\mathrm{Box}_1,\mathrm{Box}_2$ & certified squares of half-width $10^{-6}$, \S\ref{sec:method} \\
\hline
\end{tabular}
\end{center}

\section{Background}
\label{sec:background}

\subsection{Hecke triangle groups}
\label{sec:hecke}

For integer $q\ge 3$, the group $\Gq$ is a cofinite discrete subgroup
of $\mathrm{PSL}(2,\mathbb R)$ with signature $(0;2,q,\infty)$. It is
arithmetic if and only if $q\in\{3,4,6\}$ \cite{Takeuchi1977}. The
surface $X_q$ has one cusp, so the scattering matrix of the trivial
character is $1\times 1$; its determinant is a scalar meromorphic
function $\phiq$. We work throughout with the conjugated
normalization in which the cusp has width $1$, as in Hejhal
\cite[Ch.~7]{Hejhal1983}. Changing the width multiplies $\phiq$ by a
zero-free factor of the form $c^{1-2s}$, which preserves zeros, poles,
and the functional equation $\phiq(s)\phiq(1-s)=1$.

\subsection{Transfer operators}
\label{sec:to}

Mayer, M\"uhlenbruch, and Str\"omberg construct, for each integer
$q\ge 3$, a nuclear transfer operator on a finite direct sum of Hardy
spaces of disks, together with a $\pm$ decomposition under a
CF-reflection involution \cite{MMS2012}. For odd $q$ the reduced
plus-sector operator $L_{s,+}$ is the operator displayed by the
odd-$q$ incidence formulas of that paper (the three-row reduced
block structure). At $q=5$ one has $h_q=1$ and $\kappa_q=3$. The
heading printed above \cite[(34)]{MMS2012} reads
``$q=2h_q+3>5$'', while the general odd-$q$ incidence formulas, Lemma
6.3, Theorem 6.4, and the explicit $q=5$ functional equation of the
same paper are stated for odd $q\ge 5$. We use the general incidence
formula at $h_q=1$, and we record the heading conflict as an open
joint in Proposition~\ref{prop:joints}~(J8).

The Fredholm determinant $\detp$ is entire in $s$ on the domain
\begin{equation}
\label{eq:omega-star}
\Omega^*=\{\Re s>1/2\}\cup\{\Re s>0,\,\Im s>1\}
\end{equation}
on which the Hilbert-space realization used in the certificates and
the Banach-space realization of \cite{MMS2012} are identified by the
transport discussed as joint (J6) below. Both certified boxes of
Theorem~\ref{thm:twopins} lie in the second component of
$\Omega^*$.

M\"oller--Pohl \cite{MoellerPohl2013} and Adam--Pohl \cite{AdamPohl2020}
give closely related Fredholm-determinant representations of
$\Zq$ for every cofinite Hecke triangle group, via period functions
and slow/fast transfer operators. We compute with the MMS plus-sector
operator, which is the operator implemented by the archived engine.

\subsection{Selberg zeta function}
\label{sec:selberg-zeta}

The Selberg zeta function of a cofinite Fuchsian group is the
Euler product over primitive closed geodesics
\begin{equation}
\label{eq:selberg-product}
Z_\Gamma(s)=\prod_{[\gamma]}\prod_{k=0}^\infty
\bigl(1-N(\gamma)^{-(s+k)}\bigr),
\end{equation}
absolutely convergent for $\Re s>1$ and meromorphic on $\mathbb C$.
For $\Gamma=\Gq$ we write $\Zq$. The factorization
\eqref{eq:mms-factor} is \cite[Theorem 6.4]{MMS2012}: it identifies
$\Zq$ with a quotient of three Fredholm determinants, of which
$\det(1-K_s)$ is an explicit entire correction whose zeros form the
lattice
\begin{equation}
\label{eq:ks-lattice}
s=-n+i\pi k/a_q,\qquad n\in\mathbb N\cup\{0\},\quad k\in\mathbb Z,
\end{equation}
hence all have real part $\le 0$. Both boxes of
Theorem~\ref{thm:twopins} satisfy $\Re s\ge 0.41054273549473627>0$, so
they are disjoint from this lattice.

\cite[Theorem 4.10]{MMS2012} states that the MMS determinants
$\det(1-L_{s,\pm})$ are meromorphic with poles only at the real points
$s=(1-k)/2$. Both certified boxes have imaginary part strictly larger
than $5$, so $\det(1-L_{s,-})$ is holomorphic on a neighbourhood of
each box.

\subsection{The FJS divisor theorem}
\label{sec:fjs}

Friedman, Jorgenson, and Smajlovi\'c \cite{FJS2021} determine the
divisors of the scattering determinant and of the Selberg zeta
function for a cofinite Fuchsian group, including orbifolds,
with a finite-dimensional unitary representation. The scalar
trivial-character, one-cusp specialization is $h=1$, $\chi=1$, and
degree of singularity $k=1$. (Immediately after defining $k=\sum_j\dim
V_j$, the source prints a colliding ``$k=0$'' in a sentence that
evidently uses $k$ for automorphic weight; the divisor formulas retain
the degree-of-singularity parameter. The collision does not change the
set-level classification used here.)

\begin{theorem}[Friedman--Jorgenson--Smajlovi\'c, \S2.4--2.5]
\label{thm:fjs}
Let $\Gamma$ be cofinite, possibly with elliptic elements, with one
cusp and trivial character, and let $\varphi$ be its scalar scattering
determinant. Then $\varphi$ continues meromorphically to $\mathbb C$
and satisfies $\varphi(s)\varphi(1-s)=1$. The nonreal zeros of
$Z_\Gamma$ in the strip $0<\Re s<1/2$ are exactly the points
$1-\rho$, $1-\bar\rho$ where $\rho$ is a zero of $\varphi$ with
$\Re\rho>1/2$ and $\Im\rho>0$, with matching multiplicity. Every
identity, elliptic, and parabolic trivial divisor of $Z_\Gamma$ is
real, as is every discrete-spectrum zero.
\end{theorem}

The statement is the $h=1$, $k=1$ case of \cite[\S2.5, items
1--7]{FJS2021}, following Venkov \cite{Venkov1990} and Hejhal
\cite[p.~499]{Hejhal1983}. Item 6 of that list is the only source of
nonreal off-line zeros of $Z_\Gamma$. The gamma and $G_1$ factors in
the completed zeta $Z_+(s)=Z(s)/(G_1(s)\Gamma(s-1/2)^k)$ have real
divisors and meet neither box of Theorem~\ref{thm:twopins}.

Jorgenson--Smajlovi\'c \cite{JorgensonSmajlovic2017} supply adjacent
cofinite-orbifold context for the Selberg zeta derivative; we do not
consume that paper in any certificate.

\section{Certification method}
\label{sec:method}

The certificates are computer-assisted theorems in Arb ball arithmetic
\cite{Johansson2017}, accessed through python-flint, at $384$ bits.
Directed rounding is the Arb midpoint-radius enclosure: every
recorded interval is an outward rounding of the true value of the
indicated quantity. We describe exactly what the receipts certify.

\subsection{Boxes and contours}

Each pin is a closed coordinate square of half-width $10^{-6}$. Write
$\mathrm{Box}_1$ for the square centred at
\[
0.4538951800749447+5.7635372417301305\,i
\]
and $\mathrm{Box}_2$ for the square centred at
\[
0.41054373549473627+7.81976824701551188\,i.
\]
The boundary $\partial\mathrm{Box}_i$ is covered by $192$ base
straight arcs ($48$ per edge), oriented counterclockwise. Adaptive
bisection, to maximum depth $8$, refines any base arc on which the
enclosure tests below fail. The first pin uses finite-section rank
$N=160$ and accepts $284$ closed subarcs. The second pin uses
$N=288$ and accepts $452$ closed subarcs across sixteen chunks of
twelve base arcs, merged in
\texttt{chunk\_receipts/}\allowbreak\texttt{S2\_MERGED\_}\allowbreak\texttt{CONTOUR\_RECEIPT.json}.

\subsection{Finite-section determinant}

Let $P_N$ be the orthogonal projection onto the span of the first $N$
basis vectors in each Hardy component, and write
$D_N(s)=\det(I-P_N L_{s,+} P_N)$ for the finite-section determinant.
The restriction identity and the Sylvester identity give
$D_N(s)=\det_H(I-L_{s,+}P_N)$. On each closed subarc $A$ with
midpoint $s_0$ and radius $r$, Arb inversion of $I-M(s)$ over $A$
produces a bound $H\ge \sup_A\bigl|\mathrm{tr}\bigl((I-M(s))^{-1}
M'(s)\bigr)\bigr|$. Jacobi's formula then yields $|D_N'(s)|\le H
|D_N(s)|$. Writing $rH$ for the product of $r$ and this $H$, the
gate $rH<1$ produces an explicit radius $G$ such that the disk
$D_N(s_0)+\overline{B}(0,rG)$ contains $D_N(s)$ for every $s\in A$.
This Taylor enclosure is recorded per leaf.

\subsection{Tail, inflation, and homotopy}

A Gohberg--Krein/Simon determinant comparison
\cite[Theorem 4.2]{Simon1977} supplies a uniform bound
\[
\bigl|\det_H(I-L_{s,+})-D_N(s)\bigr|\le F_R
:=T_{\mathrm{tail}}(N)\,\exp\bigl(1+2B_{\mathrm{same}}(N)\bigr)
\]
on the whole closed box, once column-norm bounds for $L_{s,+}$ and
$L_{s,+}P_N$ and an input-tail bound $T_{\mathrm{tail}}(N)$ are
enclosed. The receipts record these constants. For the second pin,
$F_R\le 2.0894485\cdot 10^{-8}$ at $N=288$. Each leaf also records
the $F_R$-inflated enclosure obtained by expanding the Taylor disk
by $F_R$. The gate that this inflated disk excludes $0$ certifies
that the straight-line homotopy from $D_N$ to $\det_H(I-L_{s,+})$
never vanishes on $A$. Consequently the winding of $D_N$ along
$\partial\mathrm{Box}_i$ equals the winding of the Fredholm
determinant.

\subsection{Winding}

Argument increments are taken from a convex overlap-midpoint polygon
homotopic to the image curve inside the nonzero inflated disks.
The merged second-pin winding ball is
\[
[0.999\ldots996472261\pm 1.22\cdot 10^{-113}],
\]
which pins the integer winding to $1$. The first-pin winding ball
likewise pins to $1$. By the argument principle, $\detp$ has
exactly one zero, counted with multiplicity, in each box. A zero of
total multiplicity one is simple.

\subsection{What the receipts do not certify}

The receipts do not evaluate $\Zf$ or $\phif$. They do not certify
that $\det(1-L_{s,-})$ is nonvanishing in either box. They do not
by themselves identify the Hilbert-space determinant with the MMS
Banach determinant, nor do they resolve the printed $q>5$ heading
above \cite[(34)]{MMS2012}. Those steps are the open or cited joints
of \S\ref{sec:joints}.

Figure~\ref{fig:qop} records the finite-section operator-norm bound
$q_{\mathrm{Op}}^{\mathrm{upper}}$ on $1024$ depth-$8$
closed-contour checker leaves, four arcs of $256$ leaves. Every
value lies in $[0.2188,0.3273]$, strictly below the square-box
predictor $1/(1+\sqrt{2})\approx 0.41421$. The four per-arc
histograms overplot.

\begin{figure}[t]
\centering
\includegraphics[width=0.85\linewidth]{fig1_qop_hist.pdf}
\caption{Per-arc distribution of $q_{\mathrm{Op}}^{\mathrm{upper}}$
across $1024$ depth-$8$ closed-contour checker leaves, arcs $0$--$3$,
with the square-box predictor $1/(1+\sqrt{2})\approx 0.41421$ marked.
The recorded range is $[0.2188,0.3273]$. The four series overplot.}
\label{fig:qop}
\end{figure}

\section{Two certified zeros of $\det(1-L_{s,+})$}
\label{sec:theorem}

\begin{theorem}[computer-assisted]
\label{thm:twopins}
Let $L_{s,+}$ be the plus-sector transfer operator of $G_5$ in the
sense of \S\ref{sec:to}. The Fredholm determinant $\detp$ has zeros
$s_1$ and $s_2$ of multiplicity one, with
\begin{gather*}
\Re s_1\in[0.45389418007494470,\,0.45389618007494470],\\
\Im s_1\in[5.7635362417301305,\,5.7635382417301305];
\end{gather*}
\begin{gather*}
\Re s_2\in[0.41054273549473627,\,0.41054473549473627],\\
\Im s_2\in[7.81976724701551188,\,7.81976924701551188].
\end{gather*}
The closed real-part intervals are disjoint, with separation
$0.04334944458020843$.
\end{theorem}

\begin{proof}
The first-pin receipt
\texttt{certificates/pin1\_flagship/}\allowbreak\texttt{R3B\_FLAGSHIP\_}\allowbreak\texttt{CERT\_RECEIPT.json}
certifies winding $1$ for $D_{160}$ on $\partial\mathrm{Box}_1$, with
every finite Taylor enclosure and every $F_R$-inflated enclosure
excluding zero, minimum certified margin (finite lower bound minus
$F_R$) at least $3.43786\cdot 10^{-8}$, and $rH\le 0.359<1$. The
homotopy of \S\ref{sec:method} transfers the winding to
$\det_H(I-L_{s,+})$. The argument principle therefore places exactly
one zero of $\detp$, counted with multiplicity, in $\mathrm{Box}_1$.
The second-pin merged receipt
\texttt{chunk\_receipts/}\allowbreak\texttt{S2\_MERGED\_}\allowbreak\texttt{CONTOUR\_RECEIPT.json}
does the same for $N=288$ on $\partial\mathrm{Box}_2$: sixteen chunks
cover the $192$ base arcs, $452$ accepted leaves, merged winding $1$,
minimum certified margin at least $3.06455\cdot 10^{-8}$, and
$F_R\le 2.0894485\cdot 10^{-8}$. Multiplicity one is the winding.
The displayed endpoints are the coordinate bounds of the two boxes.
The closed-interval distance between the real-part projections is
the exact decimal
$0.45389418007494470-0.41054473549473627=0.04334944458020843$,
re-verified with Arb directed rounding.
\end{proof}

Winding $1$ certifies a simple zero of $\detp$ only. If
\eqref{eq:mms-factor} applies at $q=5$ and $\det(1-K_s)$ is
nonvanishing on the box, then $\Zf$ vanishes at $s_i$ with
multiplicity at least one. Simplicity of the $\Zf$-zero would
additionally require $\det(1-L_{s,-})$ nonvanishing at $s_i$, which
is not certified.

\begin{remark}
An $N=128$ control on $\mathrm{Box}_2$ fails the self-consistency
gate ($rH>1$ on a base arc). The $N=288$ certificate is therefore
not an artefact of a uniformly loose tail.
\end{remark}

\section{Dependency proposition}
\label{sec:joints}

Theorem~\ref{thm:twopins} is unconditional on the certificates of
\S\ref{sec:method}. The scattering interpretation --- two zeros of
$\phif$ with distinct real parts in $(1/2,1)$, hence no single
vertical line $\Re s=c$ containing every nonreal strip zero of
$\phif$ --- uses further joints. Each item below is marked
\textsc{certified} (an Arb receipt in the archive),
\textsc{cited} (a published theorem, with hypotheses checked
in-line), or \textsc{open}.

\begin{proposition}
\label{prop:joints}
The passage from Theorem~\ref{thm:twopins} to two zeros of $\phif$
with distinct real parts in $(1/2,1)$ factors as follows.
\begin{enumerate}
\item[\textup{(J1)}]
Exact $q=5$ plus-sector operator and determinant normalization:
the operator of \S\ref{sec:to}, sign $+1$, head split $4$, Hardy
direct sum, as implemented by the archived evaluation engine consumed by
both receipts.
\textsc{certified}.

\item[\textup{(J2)}]
Whole-contour Arb enclosure of $D_N$ and of the $F_R$-inflated
disks, with zero excluded on every accepted subarc of
$\partial\mathrm{Box}_i$.
\textsc{certified}.

\item[\textup{(J3)}]
Winding equal to $1$ on each contour. This certifies a simple zero
of $\detp$ only.
\textsc{certified}.

\item[\textup{(J4)}]
Finite section to the infinite Hilbert-space determinant: the
identity $D_N=\det_H(I-L_{s,+}P_N)$ together with the
Gohberg--Krein/Simon bound and the enclosed constants
$T_{\mathrm{tail}}(N)$, $B_{\mathrm{same}}(N)$, $F_R$.
\textsc{cited}+\textsc{certified}.

\item[\textup{(J5)}]
Straight-line homotopy and the argument principle on
$\partial\mathrm{Box}_i$, using (J2) and (J4).
\textsc{cited}+\textsc{certified}.

\item[\textup{(J6)}]
Equality of the Hilbert-space determinant and the MMS Banach
determinant on $\Omega^*$ in \eqref{eq:omega-star}
(``Hilbert-to-Banach transport''). Both boxes lie in
$\{\Re s>0,\,\Im s>1\}$. The identification is recorded as a
paper-proof in the archive and is not reproduced here.
\textsc{open}.

\item[\textup{(J7)}]
Exclusion of every zero of $\det(1-K_s)$ from each closed box. By
\eqref{eq:ks-lattice} those zeros have real part $\le 0$, while
both boxes satisfy $\Re s>0$.
\textsc{cited}.

\item[\textup{(J8)}]
The factorization \eqref{eq:mms-factor} at $q=5$, i.e.\ that
\cite[Theorem 6.4]{MMS2012} applies despite the printed heading
``$q=2h_q+3>5$'' above \cite[(34)]{MMS2012}. The general odd-$q$
incidence formulas and the explicit $q=5$ functional equation of
that paper support the identification; the heading conflict is not
dissolved in print.
\textsc{open}.

\item[\textup{(J9)}]
A zero of $\detp$ is a zero of $\Zf$ of multiplicity at least that
of the plus factor. This follows from \eqref{eq:mms-factor} once
(J7) and (J8) hold, using holomorphy of $\det(1-L_{s,-})$ near
each box \cite[Theorem 4.10]{MMS2012}. Simplicity of the
$\Zf$-zero is not claimed.
\textsc{cited}, conditional on \textup{(J8)}.

\item[\textup{(J10)}]
A nonreal left-strip zero of $\Zf$ is a pole of $\phif$ of the
same order, with no completed, trivial, or gamma divisor in either
box: Theorem~\ref{thm:fjs}, item 6 of \cite[\S2.5]{FJS2021},
one-cusp scalar specialization.
\textsc{cited}.

\item[\textup{(J11)}]
The functional equation $\phif(s)\phif(1-s)=1$ reflects a pole of
order $m$ at $s_i$ to a zero of order $m$ at $1-s_i$. No
conjugation enters.
\textsc{cited}.

\item[\textup{(J12)}]
Transport of the two rectangles under $s\mapsto 1-s$ and
distinctness of the real parts. The image intervals are
\begin{gather*}
\Re\rho_1\in[0.54610381992505530,\,0.54610581992505530],\\
\Im\rho_1\in[-5.7635382417301305,\,-5.7635362417301305];
\end{gather*}
\begin{gather*}
\Re\rho_2\in[0.58945526450526373,\,0.58945726450526373],\\
\Im\rho_2\in[-7.81976924701551188,\,-7.81976724701551188],
\end{gather*}
with the same closed-interval separation $0.04334944458020843$.
Both real-part intervals lie in $(1/2,1)$.
\textsc{certified}.

\item[\textup{(J13)}]
Membership of $(\phif,\mathcal D_5)$ in a generic class of
scattering pairs. That membership is not used in this paper.
\textsc{open}.
\end{enumerate}
Joints \textup{(J1)}--\textup{(J5)} and the $s$-box arithmetic in
\textup{(J12)} are the support of Theorem~\ref{thm:twopins}. The
scattering and no-single-line consequence below is conditional on
every \textsc{cited} or \textsc{open} joint among
\textup{(J6)}--\textup{(J12)}, and is thereby severed from the
unconditional two-pin theorem. Joint \textup{(J13)} is not consumed.
\end{proposition}

\begin{corollary}[conditional on \textup{(J6)}--\textup{(J12)}]
\label{cor:scattering}
Assume joints \textup{(J6)}--\textup{(J11)}. Then $\Zf$ has zeros
at $s_1$ and $s_2$ of multiplicity at least one, $\phif$ has poles
at $s_i$ and zeros $\rho_i=1-s_i$ of the same order, and the
real-part intervals of $\rho_1$ and $\rho_2$ are those displayed in
\textup{(J12)}. In particular, the nonreal zeros of $\phif$ in
$1/2<\Re s<1$ lie on no single vertical line $\Re s=c$.
\end{corollary}

The last clause is elementary once two zeros with distinct real
parts are granted: for any $c\in\mathbb R$, at least one of the two
closed real-part intervals excludes $c$. Joint (J8) remains open,
so Corollary~\ref{cor:scattering} is not a theorem of this paper.

\section{Related work}
\label{sec:related}

The novelty claim is narrow: these are, to our knowledge, the first
\emph{interval-certified} localizations of zeros of the $G_5$
plus-sector transfer-operator determinant, as against the heuristic
numerics of \cite{Stromberg2009}. The analytic apparatus is
classical. Table~\ref{tab:delta} records, for each load-bearing
predecessor, what it proves and what the present paper adds.

\begin{table}[t]
\centering
\small
\caption{Theorem delta. ``Adds'' means relative to the cited result,
not a claim of priority against the rest of the table.}
\label{tab:delta}
\begin{tabular}{@{}>{\raggedright\arraybackslash}p{0.20\textwidth}
>{\raggedright\arraybackslash}p{0.38\textwidth}
>{\raggedright\arraybackslash}p{0.32\textwidth}@{}}
\hline
Source & Proves & This paper \\
\hline
Selberg \cite{Selberg1990} &
Weighted distribution of poles of Eisenstein series on hyperbolic
surfaces; Jensen/Littlewood counting engine. &
Not used. No counting theorem is claimed. \\
Hejhal \cite{Hejhal1983} &
Thm.\ 7.11 / Cor.\ 7.12: for all sufficiently large $N$, zeros
(resp.\ poles) of $\varphi_N$ in any prescribed rectangle touching
$\Re s=1/2$. Memoirs 469 \cite{Hejhal1992}: numerical Hecke spectra.
&
Certificates at the single group $G_5$, with explicit boxes; no
large-$q$ statement. \\
Kelmer \cite{Kelmer2015} &
Higher-dimensional distribution of Eisenstein poles; surface case
attributed to Selberg. Thm.\ 3 and \S4. &
Not used. \\
Phillips--Sarnak \cite{PhillipsSarnak1985,PhillipsSarnak1992,PhillipsSarnakHecke} &
Deformation of the augmented spectrum; cusp-form dissolution into
resonances under a Fermi golden rule; spectrum of Hecke groups. &
No deformation is used. $G_5$ is a fixed orbifold. \\
Garbin--Jorgenson \cite{GarbinJorgenson2018} &
Quantitative elliptic degeneration of the Hecke family; reproduces
Hejhal 7.11/7.12; spectral counting $c_0(T)\log Q+O((\log Q)^{3/4})$.
&
No $q\to\infty$ statement. \\
Str\"omberg \cite{Stromberg2009} &
Numerical zeros of $\Zq$ for $q=3,4,5$ from truncated transfer
operators; explicitly non-rigorous. &
Interval certificates for two $G_5$ zeros of $\detp$, winding $1$. \\
MMS \cite{MMS2012} &
Nuclear transfer operator of $\Gq$; factorization
\eqref{eq:mms-factor}; meromorphy of the sector determinants. &
Evaluates $\detp$ in ball arithmetic at $q=5$. The $q>5$ heading
is left open (J8). \\
M\"oller--Pohl \cite{MoellerPohl2013} &
Period functions for every cofinite $\Gq$; Selberg zeta as a
Fredholm determinant of accelerated transfer operators. &
Different operator; not evaluated here. \\
Adam--Pohl \cite{AdamPohl2020} &
Laplace eigenfunctions versus zeros of $\Zq$ via slow/fast
transfer operators, finite and infinite area. &
Not evaluated here. \\
FJS \cite{FJS2021} &
Divisor of $\varphi$ and of $Z_\Gamma$ for cofinite orbifolds
with unitary representations (Theorem~\ref{thm:fjs}). &
Cited for the conditional promotion (J10), not re-proved. \\
Jorgenson--Smajlovi\'c \cite{JorgensonSmajlovic2017} &
Zeros of the derivative of the Selberg zeta on finite-volume
surfaces. &
Context only; not consumed. \\
\hline
\end{tabular}
\end{table}

Fried \cite{Fried1986} is the historical source of dynamical zeta
functions as Fredholm determinants; it is not a finite-area Hecke
cusp theorem. Borthwick--Judge--Perry \cite{BJP2005} identify
resonances with Selberg-zeta divisors for geometrically finite
surfaces, typically in infinite area.

\paragraph{An open question.}
Does the fixed-$q$ statement ``for every integer $q\ge 3$, the
scattering determinant $\phiq$ has infinitely many nonreal zeros in
$\Re s>1/2$'' follow as an immediate corollary of Selberg
\cite{Selberg1990}, Hejhal \cite{Hejhal1983}, and Kelmer
\cite{Kelmer2015} in orbifold generality? Hejhal's Theorem 7.11
gives zeros in prescribed rectangles for all sufficiently large
$q$, not for each fixed $q$. Kelmer's global setup is torsion-free.
We do not claim the fixed-$q$ growth law as a theorem, and we do
not claim to have settled this comparison.

\section{Machine verification}
\label{sec:lean}

The Lean~4 artifacts in the archive formalize elementary
conditional implications between stated hypotheses. No analytic or
spectral content is formalized: there is no Lean definition of
$\phif$, of $L_{s,+}$, of $\Zf$, or of a Hecke triangle group, and
no theorem of MMS or FJS is in Lean.

The file \texttt{lean/\allowbreak v34\_TwoPinNoLine/\allowbreak TwoPinNoLine.lean} proves
four statements: an abstract set of complex numbers containing two
nonreal right-strip members with distinct real parts admits no
vertical line through all such members; the exact rational gap
$4334944458020843/10^{17}$ between the two pin intervals; disjointness
of those intervals; and distinctness of real parts for any pair of
reals lying in the two intervals. The zeros are members of an
abstract set. Multiplicity is not present: the formalization does
not distinguish a simple zero from a multiple zero, and it does not
treat the Selberg-zeta multiplicity lower bound of (J9).

The file
\texttt{lean/\allowbreak v33\_Scat1Lemma31Reflection/\allowbreak }\allowbreak\texttt{Scat1Lemma31Reflection.lean}
proves that if a meromorphic function on an open set stable under
$s\mapsto 1-s$ satisfies $\varphi(s)\varphi(1-s)=1$ away from poles,
then a pole of meromorphic order $m\ge 1$ at $s_*$ implies a zero of
meromorphic order $m$ at $1-s_*$. Multiplicity is the mathlib
meromorphic order (a pole of order $m$ is order $-m$). The
hypotheses are abstract; $\phif$ is not instantiated.

Both files are sorry-free and depend only on the standard classical
axioms \texttt{propext}, \texttt{Classical.choice}, and
\texttt{Quot.sound}. They are regression checks for the elementary
steps in (J11) and (J12), not independent evidence for
Theorem~\ref{thm:twopins}.

\section{Data availability}
\label{sec:data}

Supporting material will be deposited at [DOI pending]. Until that
deposit, the archive consists of the files named below, with paths
relative to the archive root. No statement of archival completeness
is made in advance of the DOI.

\begin{sloppypar}\raggedright
\begin{itemize}
\item
\texttt{certificates/pin1\_flagship/}\allowbreak\texttt{R3B\_FLAGSHIP\_}\allowbreak\texttt{CERT\_RECEIPT.json}
and \texttt{R3B\_FLAGSHIP\_CERT.md}: first-pin closed-contour
certificate, $N=160$, winding $1$.
\item
\texttt{chunk\_receipts/}\allowbreak\texttt{S2\_MERGED\_}\allowbreak\texttt{CONTOUR\_RECEIPT.json}: merged
second-pin closed-contour certificate, $N=288$,
\texttt{merged\_winding=1}, together with the sixteen chunk
receipts \texttt{S2\_CHUNK\_a000-012.json}, \ldots,
\texttt{S2\_CHUNK\_a180-192.json}.
\item
\texttt{S2\_MERGED\_REFEREE\_FABLE.md} and
\texttt{S2\_MERGED\_REFEREE\_SOL.md}: the two merged-contour referee
reports.
\item
\texttt{lean/\allowbreak v34\_TwoPinNoLine/\allowbreak TwoPinNoLine.lean} and
\texttt{lean/\allowbreak v33\_Scat1Lemma31Reflection/\allowbreak }\allowbreak\texttt{Scat1Lemma31Reflection.lean},
with \texttt{lean/VERIFICATION.md}.
\item
\texttt{fig1\_qop\_hist.pdf}, reproduced from
\texttt{shard\_receipts\_d8/} by \texttt{make\_fig1.py}.
\end{itemize}
\end{sloppypar}

\begin{thebibliography}{99}

\bibitem{AdamPohl2020}
A.~Adam and A.~D.~Pohl,
\emph{A transfer-operator-based relation between Laplace
eigenfunctions and zeros of Selberg zeta functions},
Ergodic Theory Dynam.\ Systems,
DOI 10.1017/etds.2018.51;
arXiv:1606.09109.

\bibitem{BJP2005}
D.~Borthwick, C.~Judge, and P.~A.~Perry,
\emph{Selberg's zeta function and the spectral geometry of
geometrically finite hyperbolic surfaces},
Comment.\ Math.\ Helv.\ \textbf{80} (2005), 483--515.
DOI 10.4171/CMH/23.

\bibitem{FJS2021}
J.~S.~Friedman, J.~Jorgenson, and L.~Smajlovi\'c,
\emph{Super-zeta functions and regularized determinants associated
to cofinite Fuchsian groups with finite-dimensional unitary
representations},
Lett.\ Math.\ Phys.\ \textbf{111} (2021), article 15;
arXiv:2011.12795.

\bibitem{Fried1986}
D.~Fried,
\emph{The zeta functions of Ruelle and Selberg. I},
Ann.\ Sci.\ \'Ecole Norm.\ Sup.\ (4) \textbf{19} (1986), 491--527.
DOI 10.24033/asens.1515.

\bibitem{GarbinJorgenson2018}
D.~Garbin and J.~Jorgenson,
\emph{Spectral asymptotics on sequences of elliptically degenerating
Riemann surfaces},
Enseign.\ Math.\ \textbf{64} (2018), 161--206.
DOI 10.4171/LEM/64-1/2-7.

\bibitem{Hejhal1983}
D.~A.~Hejhal,
\emph{The Selberg Trace Formula for $\mathrm{PSL}(2,\mathbb R)$,
Vol.\ 2},
Lecture Notes in Mathematics 1001, Springer, 1983.
DOI 10.1007/BFb0061302.
(Thm.\ 7.11 / Cor.\ 7.12; divisor discussion p.~499.)

\bibitem{Hejhal1992}
D.~A.~Hejhal,
\emph{Eigenvalues of the Laplacian for Hecke triangle groups},
Mem.\ Amer.\ Math.\ Soc.\ \textbf{97} (1992), no.~469.
DOI 10.1090/memo/0469.

\bibitem{Johansson2017}
F.~Johansson,
\emph{Arb: efficient arbitrary-precision midpoint-radius interval
arithmetic},
IEEE Trans.\ Comput.\ \textbf{66} (2017), no.~8, 1281--1294.
DOI 10.1109/TC.2017.2690633.

\bibitem{JorgensonSmajlovic2017}
J.~Jorgenson and L.~Smajlovi\'c,
\emph{On the distribution of zeros of the derivative of the Selberg
zeta-function associated to finite volume Riemann surfaces},
Nagoya Math.\ J.\ \textbf{226} (2017);
DOI 10.1017/nmj.2016.52.

\bibitem{Kelmer2015}
D.~Kelmer,
\emph{On the distribution of poles of Eisenstein series and the
length spectrum of hyperbolic manifolds},
Int.\ Math.\ Res.\ Not.\ IMRN (2015).
DOI 10.1093/imrn/rnv051;
arXiv:1402.4780.

\bibitem{MMS2012}
D.~Mayer, T.~M\"uhlenbruch, and F.~Str\"omberg,
\emph{The transfer operator for the Hecke triangle groups},
Discrete Contin.\ Dyn.\ Syst.\ \textbf{32} (2012), 2453--2484.
DOI 10.3934/dcds.2012.32.2453;
arXiv:0912.2236.

\bibitem{MoellerPohl2013}
M.~M\"oller and A.~D.~Pohl,
\emph{Period functions for Hecke triangle groups, and the Selberg
zeta function as a Fredholm determinant},
Ergodic Theory Dynam.\ Systems \textbf{33} (2013), 247--283.
DOI 10.1017/S0143385711000794.

\bibitem{PhillipsSarnak1985}
R.~S.~Phillips and P.~Sarnak,
\emph{On cusp forms for co-finite subgroups of
$\mathrm{PSL}(2,\mathbb R)$},
Invent.\ Math.\ \textbf{80} (1985), 339--364.
DOI 10.1007/BF01388610.

\bibitem{PhillipsSarnakHecke}
R.~S.~Phillips and P.~Sarnak,
\emph{On the spectrum of the Hecke groups},
Duke Math.\ J.\ \textbf{52} (1985), 211--221.
DOI 10.1215/S0012-7094-85-05212-3.

\bibitem{PhillipsSarnak1992}
R.~S.~Phillips and P.~Sarnak,
\emph{Perturbation theory for the Laplacian on automorphic
functions},
J.\ Amer.\ Math.\ Soc.\ \textbf{5} (1992), 1--32.
DOI 10.1090/S0894-0347-1992-1127079-8.

\bibitem{Selberg1990}
A.~Selberg,
\emph{Remarks on the distribution of poles of Eisenstein series},
in: Festschrift in honor of I.~I.~Piatetski-Shapiro,
Israel Math.\ Conf.\ Proc.\ \textbf{3} (1990), 251--278.

\bibitem{Simon1977}
B.~Simon,
\emph{Notes on infinite determinants of Hilbert space operators},
Adv.\ Math.\ \textbf{24} (1977), 244--273.
DOI 10.1016/0001-8708(77)90057-3.

\bibitem{Stromberg2009}
F.~Str\"omberg,
\emph{Computation of Selberg zeta functions on Hecke triangle
groups},
Experiment.\ Math.\ \textbf{18} (2009), 197--214;
arXiv:0804.4837.

\bibitem{Takeuchi1977}
K.~Takeuchi,
\emph{Arithmetic triangle groups},
J.\ Math.\ Soc.\ Japan \textbf{29} (1977), 91--106.
DOI 10.2969/jmsj/02910091.

\bibitem{Venkov1990}
A.~B.~Venkov,
\emph{Spectral theory of automorphic functions and its
applications},
Mathematics and its Applications (Soviet Series) 51, Kluwer,
Dordrecht, 1990.
(English translation of Trudy Mat.\ Inst.\ Steklov \textbf{153}
(1981); Thm.\ 3.5.)

\end{thebibliography}

\end{document}
